\documentclass[10pt,a4paper]{article}

\usepackage[margin=1in]{geometry}
\usepackage{amsmath,amssymb}
\usepackage{hyperref}
\usepackage{booktabs}
\usepackage{listings}
\usepackage{xcolor}

\hypersetup{colorlinks=true, linkcolor=blue!50!black, citecolor=blue!50!black, urlcolor=blue!50!black}

\title{LLL Failures on the Boneh--Venkatesan Lattice Were\\
       Floating-Point Artifacts: A Corrected Note}
\author{Anonymous}
\date{}

\newcommand{\FF}{\mathbb{F}}
\newcommand{\ZZ}{\mathbb{Z}}

\begin{document}
\maketitle
\emergencystretch=2em

\begin{abstract}
An earlier version of this note (PDF of 2026-10-11) reported two LLL
failure modes for the Boneh--Venkatesan Hidden-Number-Problem lattice
on deployed ECDSA curves: a ``curve-specific'' failure for secp256k1
at 256 bits and a uniform failure at $\geq 384$ bits attributed to an
undersized iteration cap. \textbf{Both were bugs in our floating-point
Gram--Schmidt orthogonalisation, not properties of the curves or of
LLL.} For secp256k1, and for every 384-bit curve, the squared
Gram--Schmidt norms overflow IEEE double precision ($n^4 >$
\texttt{f64::MAX}). P-521 still fails on the scaled \texttt{f64}
path, for a reason we have not isolated. Both fixes (scaling, and a
2048-bit fixed-point Gram--Schmidt) were in the companion code before
that version was built. With them, a re-run on 2026-10-16 recovers the
key for every distinct seed pair at 256 and 384 bits, secp256k1
included (3 of 3 distinct seed pairs, 6 probe runs), and on P-521 with
2048-bit fixed-point Gram--Schmidt (3 of 3; at
$k_{\text{bits}} = 384$). We withdraw the curve-specificity claim and every conclusion
drawn from it. What remains is a numerical lesson: an
$\texttt{f64}$ Gram--Schmidt step silently fails once the group order
passes about $2^{256}$ (basis entries about $2^{512}$).
\end{abstract}

\section*{Correction (2026-10-16)}

This version replaces the note of 2026-10-11. An integrity audit of
our results (the \texttt{sheriff} case \texttt{SHF-20261016-3caa6f},
finding P1, in the companion research repository) found that the note
reported as curve properties two failures that our own analysis
(\nolinkurl{research/notes/lattice-hnp/RESEARCH_LLL_GS_ANALYSIS.md},
\S10) had traced to numerical bugs and fixed in May 2026. The
following claims of the earlier version are \textbf{withdrawn}:
\begin{itemize}
\item that secp256k1's LLL failure is ``a curve-specific arithmetic
      phenomenon'', ``empirically robust'' and ``reproducible'';
\item that the $\geq 384$-bit failures were caused by the LLL iteration
      cap;
\item that LLL-convergence benchmarks are curve-specific and that
      reported HNP attack costs do not transfer between curves;
\item the proposed curve-selection criterion against
      near-power-of-two group orders, and the open questions built on
      the curve-specificity claim.
\end{itemize}
The pre-fix measurements are kept below (Section~\ref{sec:prefix}) as
the record of what the buggy code did.

\section*{Correction (2026-10-17)}

A re-audit (\texttt{sheriff} case \texttt{SHF-20261017-e4d7f6},
findings P-N1 and P-N2) found two defects in the 2026-10-16 version.
(i) It asserted a mechanism for the P-521 failure on the scaled
\texttt{f64} path: ``catastrophic cancellation'', in which
orthogonalising a row ``subtracts nearly equal 53-bit quantities, and
the residual swamps the true component''. That mechanism is
unsupported. With our scaling, the P-521 diagonal $n^2/2^{542}$ rounds
to exactly $2^{500}$ in \texttt{f64}, so the subtraction described is
exact on the initial basis, and the residual-dominates condition also
holds for P-256, brainpoolP256r1 and brainpoolP384r1, which recover.
The sentence is withdrawn; the measured result (scaled \texttt{f64}
0 of 3 distinct seed pairs at P-521, 2048-bit fixed point 3 of 3)
stands. (ii) Table~\ref{tab:rerun} counted probe runs of deterministic,
overlapping test functions as if they were distinct instances, and
said $k_{\text{bits}} = 3n_{\text{bits}}/4$ throughout, whereas every
P-521 probe ran at $k_{\text{bits}} = 384$ ($3/4$ of 512, not of 521).
The table now gives distinct seed pairs, runs and $k_{\text{bits}}$.

\section{Setting and notation}

The Boneh--Venkatesan lattice attack on biased-nonce
ECDSA~\cite{boneh1996,nguyen2002} recovers the secret key $d$ from
$m$ signatures $(r_i, s_i, z_i)$ in which the nonce $k_i$ is known
to satisfy $0 \leq k_i < 2^{k_{\text{bits}}}$ for some
$k_{\text{bits}} < n_{\text{bits}}$. With $a_i = s_i^{-1} r_i \bmod n$
and $t_i = s_i^{-1} z_i \bmod n$, each signature contributes the
constraint
\[
   d \cdot a_i + t_i \equiv k_i \pmod{n}.
\]

The standard $(m+2)$-dimensional integer lattice basis (rows)
\[
   B \;=\;
   \begin{pmatrix}
     n^2 \cdot I_m   & 0_{m \times 2}  \\[2pt]
     n \cdot \mathbf{a}^\top  & 2^{k_{\text{bits}}}  & 0 \\[2pt]
     n \cdot \mathbf{t}^\top  & 0  & n \cdot 2^{k_{\text{bits}}}
   \end{pmatrix}
\]
is fed to LLL with Lov\'asz parameter $\delta = 0.75$; the short
vector after reduction contains $-d \cdot 2^{k_{\text{bits}}}$ in
the $(m+1)$-st coordinate (modulo sign).

\section{What the earlier version measured}
\label{sec:prefix}

These are the pre-fix measurements as the 2026-10-11 version reported
them. They were produced by unscaled \texttt{f64} Gram--Schmidt and
are \textbf{superseded} by Section~\ref{sec:rerun}.

\begin{table}[h]
\centering
\begin{tabular}{lrrr}
\toprule
$(d_{\text{seed}}, k_{\text{seed}})$ & P-256 outcome & P-256 time & secp256k1 outcome \\
\midrule
\texttt{(0xC0FFEE, 0xC0FFEE)}           & recovered & 690~ms  & iteration cap @ 9993~ms \\
\texttt{(0xDEADBEEF, 0xBADCAFE)}        & recovered & 697~ms  & iteration cap @ 9845~ms \\
\texttt{(0x12345678, 0x9ABCDEF0)}       & recovered & 646~ms  & iteration cap @ 10119~ms \\
\bottomrule
\end{tabular}
\caption{Pre-fix (superseded): $k_{\text{bits}} = 192$, $m = 8$,
$\delta = 0.75$. Reported then as P-256 3/3, secp256k1 0/3.}
\label{tab:prefix}
\end{table}

The pre-fix sweep at $k_{\text{bits}} = 3 n_{\text{bits}}/4$, $m = 8$,
one seed, reported secp192k1, secp224k1, P-256 and brainpoolP256r1
recovering, and secp256k1, P-384, brainpoolP384r1 and P-521 hitting
the iteration cap.

\section{The causes}

The basis has rows of size about $n^2$, so the squared Gram--Schmidt
norms the reduction compares are of size about $n^4$.

\paragraph{Failure A: overflow.} IEEE double precision tops out just
below $2^{1024}$. Among the three 256-bit curves, only secp256k1's
order is close enough to $2^{256}$ that $n^4$ exceeds
\texttt{f64::MAX}; P-256's and brainpoolP256r1's orders are smaller,
and their $n^4$ stays below it (an exact rational comparison of the
standard orders~\cite{fips186,sec2}, recorded in the audit as
\texttt{checks/f64-overflow.txt}). The overflowed norms became
infinite and then \texttt{NaN}, and LLL ran to its iteration cap. The
``curve specificity'' was the position of one group order relative to
a floating-point limit. secp192k1 and secp224k1 escaped because their
$n^4$ is far below the limit, which is also why the earlier
near-power-of-two hypothesis seemed refuted. Scaling the basis by
$2^{-(\text{max\_bits} - 500)}$ before orthogonalising removes the
overflow (\texttt{src/cryptanalysis/lattice.rs}).

\paragraph{Failure B: overflow, and an unexplained P-521 failure.} At
384 bits, $n^4$ overflows for every curve, which explains the
``uniform'' failure there; the same scaling fixes P-384 and
brainpoolP384r1, so the iteration cap was not the cause. P-521 still
fails after scaling (0 of 3 distinct seed pairs), and a 2048-bit
fixed-point Gram--Schmidt (\texttt{lll\_reduce\_hp}) recovers it
(3 of 3). We have not isolated where the scaled \texttt{f64} path
loses precision at P-521; the cancellation argument of \S10.2 of the
analysis does not apply to the initial basis (correction of
2026-10-17 above), and locating the failing step would need the
Gram--Schmidt instrumented during the probe.

\section{Re-run with the fixes (2026-10-16)}
\label{sec:rerun}

We re-ran the committed probes (\nolinkurl{tests/lll_degeneracy_probe.rs})
on the fixed code. Commands, logs, parsed results and the host are in
\nolinkurl{research/lll_degeneracy_rerun_20261016/}.

\begin{table}[h]
\centering
\begin{tabular}{llrrr}
\toprule
Curve & Gram--Schmidt & $k_{\text{bits}}$ & Recovered / distinct & Runs \\
\midrule
secp192k1        & scaled \texttt{f64}  & 144 & 2 / 2 & 3 / 3 \\
secp224k1        & scaled \texttt{f64}  & 168 & 2 / 2 & 3 / 3 \\
secp256k1        & scaled \texttt{f64}  & 192 & 3 / 3 & 6 / 6 \\
P-256            & scaled \texttt{f64}  & 192 & 3 / 3 & 6 / 6 \\
brainpoolP256r1  & scaled \texttt{f64}  & 192 & 1 / 1 & 1 / 1 \\
P-384            & scaled \texttt{f64}  & 288 & 3 / 3 & 4 / 4 \\
brainpoolP384r1  & scaled \texttt{f64}  & 288 & 3 / 3 & 4 / 4 \\
P-521            & scaled \texttt{f64}  & 384 & 0 / 3 & 0 / 4 \\
P-521            & 2048-bit fixed point & 384 & 3 / 3 & 3 / 3 \\
\bottomrule
\end{tabular}
\caption{Re-run, $m = 8$, $\delta = 0.75$. ``Distinct'' counts
distinct $(d_{\text{seed}}, k_{\text{seed}})$ pairs; the probes are
deterministic, and overlapping test functions repeat pairs, so
``Runs'' counts probe executions (the 2026-10-16 version reported
runs only). $k_{\text{bits}} = 3 n_{\text{bits}}/4$ except at P-521,
where the probes use $384 = 3/4 \cdot 512$, not
$\lfloor 3 \cdot 521/4 \rfloor = 390$. Each recovered probe returned
the planted key. Source: \texttt{logs/probes.log} and
\texttt{logs/p521\_hp.log} in the re-run directory (corrected
2026-10-17).}
\label{tab:rerun}
\end{table}

Every probe at 256 and 384 bits finished in 13--27~ms on an
unisolated 4-vCPU virtual machine, against the roughly 10~s and 20~s
cap hits of the earlier version. These timings are single runs, not a
benchmark.

\section{What remains}

No curve-specific LLL behaviour remains at these parameters. The
lesson is about implementations: an \texttt{f64} Gram--Schmidt step
fails silently, by overflow, once the squared norms pass $2^{1024}$,
that is once $n$ passes about $2^{256}$ and the basis entries about
$2^{512}$. At P-521 the scaled \texttt{f64} path fails for a
reason we have not isolated, while 2048-bit fixed point succeeds. An HNP
experiment that reports an LLL failure should state its arithmetic
and rule this out before attributing the failure to the curve.

This note makes no security claim about any curve. It reports a small
number of probes, one to three distinct seed pairs per curve, at one
parameter set ($m = 8$, $k_{\text{bits}} = 3 n_{\text{bits}}/4$, except
$k_{\text{bits}} = 384$ at P-521).

\begin{thebibliography}{99}

\bibitem{boneh1996} D.~Boneh, R.~Venkatesan, ``Hardness of computing
the most significant bits of secret keys in Diffie--Hellman and
related schemes'', CRYPTO 1996.

\bibitem{nguyen2002} P.~Nguyen, I.~Shparlinski, ``The insecurity of
the digital signature algorithm with partially known nonces'',
\emph{J. Cryptol.} 2002.

\bibitem{fips186} NIST FIPS 186-4, ``Digital Signature Standard'',
2013.

\bibitem{sec2} Certicom Research, ``SEC 2: Recommended Elliptic Curve
Domain Parameters'', 2010.

\end{thebibliography}

\end{document}
